%=====================================================================
% FRONT MATTER
%=====================================================================
\pagenumbering{roman}
% Front-matter figures sit outside any chapter, so give them their own prefix
% rather than letting them inherit a misleading chapter number.
\structuralfigures{F.}
\thispagestyle{empty}

\begin{center}
\vspace*{1.0cm}

{\color{secondary}\rule{0.85\textwidth}{2pt}}\\[6pt]
{\large\textsc{\color{secondary}Bachelor of Science in Information Technology}}\\[18pt]

{\Huge\bfseries\color{primary}ARTIFICIAL INTELLIGENCE}\\[10pt]
{\LARGE\bfseries\color{primary}Search, Knowledge and Reasoning}\\[16pt]

{\color{secondary}\rule{0.85\textwidth}{2pt}}\\[16pt]

{\Large\textsc{Unified Lecture Handout}}\\[6pt]
{\large\color{gray!70!black}Complete Course in One Volume $\bullet$ Fall 2026}\\[24pt]

% --- The one application domain, announced on the title page ---
\begin{tikzpicture}
  \node[rounded corners=4pt, draw=primary, line width=1.1pt, fill=primary!6,
        text=primary, align=center, font=\small\bfseries,
        minimum width=10.2cm, minimum height=1.5cm]
       {\faIcon{tasks}\\[2pt]Every example drawn from IT project management};
\end{tikzpicture}\\[10pt]
{\small\itshape\color{gray!70!black}Task dependencies and release plans, sprints and
staffing, tickets and escalations, defects and risk.}

\vspace{1.2cm}

\begin{tcolorbox}[enhanced, width=0.86\textwidth, colback=lightbg,
  colframe=primary, boxrule=0pt, leftrule=4pt, arc=2pt, top=6pt, bottom=6pt]
\small
\textbf{\color{primary}What this volume covers.} The whole of the HEC course
\emph{Artificial Intelligence} --- every topic of the HEC 2017 outline, every
topic of the HEC 2023 outline, and every learning outcome HEC 2025 attaches to
the course --- in one resequenced volume, taught from first principles to working
code.

\smallskip
The course is deliberately \emph{symbolic}: it teaches the agent that is told the
rules and must reason with them. Learning is one signpost chapter, not a unit,
because \emph{Machine Learning} owns it three semesters later. Appendix~F states
every boundary.
\end{tcolorbox}

\vspace{1.0cm}
{\small\color{gray!60!black}Six Parts $\bullet$ Twenty-Two Chapters $\bullet$ Seven Appendices}
\vspace{1cm}
\end{center}

\newpage

%=====================================================================
\chapter*{How to Use This Handout}
\addcontentsline{toc}{chapter}{How to Use This Handout}
\markboth{How to Use This Handout}{}

This is not a reference manual to be searched; it is a \term{path} to be walked.
The chapters are ordered so that every idea rests on ideas you have already met,
and the order is not the order of the HEC outline --- three deliberate departures
are set out at the end of the Course Specification.

\section*{The six parts}

\textbf{Part~I --- Foundations of Intelligent Systems} (Chapters 1--3) fixes what
an agent is, how to specify a task before building anything, and the four data
structures every later algorithm is written in. Chapter~3 is the toolkit chapter:
everything after it says ``push the successor onto the frontier'' and means
something precise.

\textbf{Part~II --- Problem Solving by Search} (Chapters 4--6) is the engine room.
Uninformed search and why it cannot scale; heuristics and $A^{*}$, which is the
single most useful algorithm in the book; and local search for the problems where
only the final state matters.

\textbf{Part~III --- Constraints, Games and Plans} (Chapters 7--9) applies that
engine three ways: to problems where the answer must satisfy rules simultaneously,
to problems where another party is choosing against you, and to problems where the
answer is a sequence of actions.

\textbf{Part~IV --- Knowledge, Logic and Reasoning} (Chapters 10--13) changes the
question from \emph{what should I do} to \emph{what follows from what I have been
told}. It ends with a working rule engine that can explain its own conclusions ---
the artefact an auditor accepts.

\textbf{Part~V --- Reasoning Under Uncertainty} (Chapters 14--16) handles the case
where what you have been told is not certain: probability and expected utility,
Bayesian networks, and fuzzy logic.

\textbf{Part~VI --- Learning, Language and Responsible AI} (Chapters 17--22) closes
the course. One signpost chapter on learning, the symbolic half of natural language
processing, the four classic systems read through the machinery you now have,
current trends, the limits, and the project.

\begin{conceptbox}
Why search comes before logic, when both published outlines list reasoning first.
Search is the part of AI you can \emph{run} on the first day, and Chapters~7, 8
and~9 are all search with extra structure --- so teaching it first means three
chapters cost almost nothing to introduce. Logic then arrives as a change of
question rather than a change of subject. The cost of this choice is that
Chapter~10 is nine weeks in; the benefit is that by then students have written
five search algorithms and know what a state space is in their hands, not just on
a slide.
\end{conceptbox}

\section*{The furniture of every chapter}

Every chapter has the same parts, in the same order. Learn to read them once.

\rowcolors{2}{white}{lightbg}
\begin{longtable}{@{}L{4.6cm}L{10.6cm}@{}}
\toprule
\rowcolor{primary!15}
\textbf{Box} & \textbf{What it is for} \\
\midrule
\endhead
\faIcon{bullseye}~\textbf{Outcomes} & What you will be able to \emph{do}. Bloom
verbs, not topics. Use it as a checklist before the examination. \\
\faIcon{link}~\textbf{Prerequisites} & Which earlier chapters this one stands on.
If these are shaky, go back; nothing here is improved by pressing on. \\
\faIcon{tags}~\textbf{Key terms} & The vocabulary introduced. Every one appears in
the glossary of Appendix~D. \\
\faIcon{book}~\textbf{Definition} & The canonical statement, with its source. Learn
these verbatim. \\
\faIcon{lightbulb}~\textbf{Concept} & The judgement call --- what an experienced
practitioner would decide here, and why. \\
\faIcon{exclamation-triangle}~\textbf{Warning} & Something that will bite you, and
usually silently. \\
\faIcon{calculator}~\textbf{Worked example} & Real arithmetic, every step shown.
Do them with a pen. \\
\faIcon{bug}~\textbf{Common pitfalls} & The mistakes that are actually made, with
the symptom each produces. \\
\faIcon{flask}~\textbf{Try it yourself} & The laboratory. Each is a runnable file
in \texttt{code/}, named for its chapter. \\
\faIcon{question-circle}~\textbf{Checkpoint} & Three questions. Attempt them before
reading the answers in Appendix~C. \\
\faIcon{clipboard-check}~\textbf{Chapter summary} & What to retain. \\
\faIcon{pen-fancy}~\textbf{Review questions} & Examination practice: multiple
choice, short answer, applied. \\
\bottomrule
\end{longtable}

\section*{How to study with it}

\begin{enumerate}[leftmargin=2.2em, itemsep=3pt, topsep=4pt]
  \item Read the outcomes first. They tell you what the chapter is \emph{for}.
  \item Read the prose once without stopping, then again with a pen on the worked
        examples. Every number in this book was computed before it was typed; you
        should be able to reproduce each one.
  \item Run the lab \emph{before} the lecture on that chapter if you can. The code
        is short and the questions it raises are the questions the lecture answers.
  \item Do the checkpoint from memory. Then read Appendix~C and argue with it.
  \item Attempt the applied review questions in pairs. They are written to be
        discussed, not looked up.
\end{enumerate}

\begin{alertbox}[On the laboratories]
Every lab in this book runs on a plain Python~3.10 installation. Chapters~10--13
and~18--21 use nothing but the standard library; the rest add NumPy. No lab needs a
GPU, a network connection or an account with anybody. That is a deliberate
property of a symbolic course, and it is worth noticing: the methods in this book
were designed when a megabyte was a lot of memory, and they still fit in one.
\end{alertbox}

\newpage

%=====================================================================
\chapter*{Course Specification}
\addcontentsline{toc}{chapter}{Course Specification}
\markboth{Course Specification}{}

This page is the course file. It states what HEC requires of a course called
\emph{Artificial Intelligence}, and shows where in this handout each requirement
is met. Three editions are in force in different programmes, so all three are
traced.

\section*{Course particulars}

\rowcolors{2}{white}{lightbg}
\begin{longtable}{@{}L{4.2cm}L{11.0cm}@{}}
\toprule
\rowcolor{primary!15}
\textbf{Item} & \textbf{Value} \\
\midrule
\endhead
Course title & Artificial Intelligence \\
Programme & BS Information Technology. Under HEC 2023 this is a \emph{computing
core} course --- compulsory in every computing degree --- taught in Semester~3;
under HEC 2017 it was a BSCS domain core absent from BSIT; under HEC 2025 it is a
major core in Semester~V \\
Credit hours & 3 (2--3): two credit hours of lectures and one credit hour of
laboratory, which is three contact hours of lab per week (HEC 2023). HEC 2017
wrote the same course as 3+1 \\
Prerequisite & Object Oriented Programming (HEC 2023). HEC 2017 required
\emph{Data Structures and Algorithms}, which is the stronger requirement;
\chref{python} supplies the data structures this course actually needs, so the
weaker prerequisite is sufficient here \\
Duration & Sixteen teaching weeks, with a midterm after Week~8 and a final
examination after Week~16 (Appendix~G) \\
Course introduction & The HEC 2023 aim, verbatim: ``An Introduction to Artificial
Intelligence and its applications towards Knowledge Based Systems \ldots\ Python
programming language will be used to explore and illustrate various issues and
techniques in Artificial Intelligence.'' \\
Set texts & S.\,J. Russell and P.\ Norvig, \emph{Artificial Intelligence: A Modern
Approach}, 3rd ed., Prentice Hall, 2010 (both editions); G.\,F. Luger and
W.\,A. Stubblefield, \emph{AI Algorithms, Data Structures, and Idioms in Prolog,
Lisp and Java}, Pearson, 2009; R.\,O. Duda, P.\,E. Hart and D.\,G. Stork,
\emph{Pattern Classification}, Wiley, 2001 (HEC 2017); D.\ Poole and
A.\ Mackworth, \emph{Artificial Intelligence: Foundations of Computational
Agents}, 2nd ed., 2017 \\
\bottomrule
\end{longtable}

\section*{Course learning outcomes}

HEC 2025 publishes five outcomes for this course and no topic outline; HEC 2017
and HEC 2023 publish outlines and expect Bloom-levelled outcomes. The seven below
are the union, and each is traced to the chapters that deliver it.

\rowcolors{2}{white}{lightbg}
\begin{longtable}{@{}L{1.3cm}L{7.6cm}L{2.1cm}L{3.8cm}@{}}
\toprule
\rowcolor{primary!15}
\textbf{CLO} & \textbf{By the end of the course the student will be able to} &
\textbf{Bloom} & \textbf{Chapters} \\
\midrule
\endhead
CLO-1 & Describe core AI concepts, including intelligent agents, search
algorithms, knowledge representation and reasoning & C2 Understand &
1, 2, 4, 5, 10, 11 \\
CLO-2 & Formulate a real problem as a search, constraint-satisfaction or planning
problem, and implement an algorithm that solves it & C3 Apply & 3--9 \\
CLO-3 & Represent knowledge formally and derive conclusions from it by
inference & C3 Apply & 10--13 \\
CLO-4 & Reason under uncertainty and select an action by expected utility &
C3 Apply & 14--16 \\
CLO-5 & Analyze AI systems and the applications of AI algorithms in real-world
scenarios & C4 Analyze & 19, 20, and every lab \\
CLO-6 & Develop and test a working AI solution of moderate size using appropriate
tools & C6 Create & 3, 22, and every lab \\
CLO-7 & Discuss the ethical considerations and the limitations of AI systems &
A3 Valuing & 21, and the audit in 22 \\
\bottomrule
\end{longtable}

\section*{Coverage of the HEC 2017 course outline}

Every topic in the published outline, in the order HEC lists it, with the chapter
that teaches it. Nothing is omitted.

\rowcolors{2}{white}{lightbg}
\begin{longtable}{@{}L{6.4cm}L{2.0cm}L{6.6cm}@{}}
\toprule
\rowcolor{primary!15}
\textbf{HEC 2017 outline topic} & \textbf{Chapter} & \textbf{Where exactly} \\
\midrule
\endhead
Introduction; basic components of AI & 1 & The agent definition; the four
classical definitions on two axes \\
Identifying AI systems & 1 & \emph{Is It AI? A Checklist} --- the five-point test,
applied to five tools in the chapter's lab \\
Branches of AI & 1 & The branch diagram, and which branch each course in the
degree owns \\
Reasoning and knowledge representation & 10 & Semantic networks, inheritance,
frames and defaults, ontologies \\
Propositional logic & 11 & Syntax, semantics, entailment, model checking, DPLL,
resolution \\
First-order logic & 12 & Quantifiers, unification with the occurs check,
skolemisation, resolution refutation \\
Problem solving by searching: uninformed & 4 & Breadth-first, depth-first,
depth-limited, iterative deepening, uniform cost \\
Problem solving by searching: informed & 5 & Greedy best-first, $A^{*}$,
admissibility, consistency, dominance, relaxation \\
Problem solving by searching: local & 6 & Hill climbing, random restart,
simulated annealing, beam search, genetic algorithms \\
Constraint satisfaction problems & 7 & Backtracking, minimum remaining values,
degree heuristic, forward checking, arc consistency \\
Adversarial search: minimax & 8 & Game trees and the minimax value \\
Adversarial search: alpha--beta pruning & 8 & Pruning, move ordering, and the
best case it achieves \\
Adversarial search: game playing & 8 & Cutting off search, evaluation functions,
the horizon effect \\
Learning: supervised & 17 & \emph{Supervised Learning} --- the perceptron worked
numerically, and linear separability \\
Learning: unsupervised & 17 & \emph{Unsupervised Learning} --- grouping without
labels, one figure, no algorithm \\
Learning: reinforcement & 17 & \emph{Reinforcement Learning} --- the
agent--reward loop, as a signpost to \emph{Machine Learning} \\
Uncertainty in AI & 14 & Probability as degree of belief, Bayes' rule,
conditional independence, expected utility \\
Fuzzy logic & 16 & Membership functions, fuzzy operators, Mamdani inference,
defuzzification \\
Recent trends in AI & 20 & Large language models, tool-using agents,
neuro-symbolic systems \\
Case study of AI systems & 19 & The General Problem Solver, Eliza, Student and
Macsyma, each reconstructed in code \\
Analysis of AI systems & 19 & The six-column framework --- task, environment,
representation, inference, knowledge source, failure mode --- applied to all
four \\
\bottomrule
\end{longtable}

\section*{Coverage of the HEC 2023 course outline}

\rowcolors{2}{white}{lightbg}
\begin{longtable}{@{}L{6.4cm}L{2.0cm}L{6.6cm}@{}}
\toprule
\rowcolor{primary!15}
\textbf{HEC 2023 outline topic} & \textbf{Chapter} & \textbf{Where exactly} \\
\midrule
\endhead
Introduction to AI and its applications towards knowledge-based systems & 1, 13 &
The field and its branches in Chapter~1; a working rule-based expert system in
Chapter~13 \\
Introduction to reasoning and knowledge representation & 10 & What a
representation must support, and four that do \\
Problem solving by searching: informed, uninformed, heuristics, local & 4, 5, 6 &
The whole of Part~II \\
Minimax algorithm, alpha--beta pruning, game playing & 8 & Chapter~8 entire \\
Case studies: General Problem Solver & 9, 19 & Means--ends analysis in
Chapter~9, where the machinery belongs; the system analysed in Chapter~19 \\
Case studies: Eliza & 18, 19 & Pattern substitution in Chapter~18; the system
analysed in Chapter~19 \\
Case studies: Student & 19 & Reconstructed using the unification of Chapter~12 \\
Case studies: Macsyma & 19 & Reconstructed as rule rewriting, using Chapter~13 \\
Learning from examples & 17 & The three paradigms, and where this course
stops \\
Artificial neural networks & 17 & \emph{Artificial Neural Networks} --- the
perceptron, the delta rule by hand, and the exclusive-or limit \\
Natural language processing & 18 & The symbolic pipeline: grammars, chart
parsing, pattern substitution, n-grams \\
Recent trends in AI and applications of AI algorithms & 20 & Chapter~20
entire \\
Python will be used throughout & 3, all labs & Chapter~3 fixes the
representations; all 22 labs are runnable Python in \texttt{code/} \\
\bottomrule
\end{longtable}

\section*{Coverage of the HEC 2025 learning outcomes}

HEC 2025 publishes outcomes only --- no prerequisite, no outline, no reading list.

\rowcolors{2}{white}{lightbg}
\begin{longtable}{@{}L{7.6cm}L{2.0cm}L{5.4cm}@{}}
\toprule
\rowcolor{primary!15}
\textbf{HEC 2025 outcome} & \textbf{Chapters} & \textbf{Met by} \\
\midrule
\endhead
Describe core AI concepts including search algorithms, knowledge representation,
and reasoning & 1, 2, 4--5, 10--12 & CLO-1 and CLO-3 \\
Implement basic AI algorithms for problem-solving and decision-making & 3--9, 14 &
CLO-2 and CLO-4. \emph{Decision-making} is met specifically by \emph{From Belief
to Decision} in Chapter~14: expected utility and the value of information \\
Analyze the applications of AI in real-world scenarios & 19, 20 & CLO-5, and the
six-column analysis framework \\
Discuss ethical considerations and limitations of AI systems & 21 & CLO-7.
\emph{Limitations} is read broadly: Chapter~21 covers computational limits and
brittleness as well as ethics \\
Develop simple AI models using appropriate tools and frameworks & 3, 22 & CLO-6.
The ``tools'' here are the Python standard library and NumPy, deliberately ---
see the note below \\
\bottomrule
\end{longtable}

\begin{alertbox}[On ``tools and frameworks'']
HEC 2025's fifth outcome asks for AI models built with appropriate tools and
frameworks. This course uses the Python standard library and NumPy and no AI
framework at all, and that is a deliberate choice that should be defended rather
than hidden.

A framework is appropriate when it saves work you understand. For a first course
in symbolic AI it does the opposite: a constraint solver called through a library
teaches nothing about arc consistency, and a planner behind an API teaches nothing
about means--ends analysis. Every algorithm in this book is short enough to write
from scratch in one laboratory session, and writing it is the learning. Students
who go on to \emph{Programming for Artificial Intelligence} and \emph{Machine
Learning} meet the frameworks there, with the understanding to judge them.
\end{alertbox}

\section*{Topics taught here that the published outlines do not name}

Four chapters go beyond the letter of both outlines. They are listed here so that
the departure is explicit, each with the clause it rests on.

\rowcolors{2}{white}{lightbg}
\begin{longtable}{@{}L{3.1cm}L{8.2cm}L{3.7cm}@{}}
\toprule
\rowcolor{primary!15}
\textbf{Chapter} & \textbf{Why it is here} & \textbf{Clause relied on} \\
\midrule
\endhead
9 Planning & Search with structure. It reuses Part~II wholesale, and it is where
the General Problem Solver --- which HEC 2023 names as a required case study ---
actually belongs & HEC 2023, ``Case Studies: General Problem Solver'' \\
10 Knowledge Representation & Semantic networks, frames and ontologies are the
representations that Chapters~11--13 reason over & HEC 2017 and 2023,
``Reasoning and Knowledge Representation'' \\
13 Rule-Based Expert Systems & The working knowledge-based system the 2023 aim
points at, and the only place explanation is built rather than discussed &
HEC 2023, ``applications towards Knowledge Based Systems'' \\
15 Bayesian Networks & The representation that makes uncertainty computable
rather than merely acknowledged & HEC 2017, ``Uncertainty in AI'' \\
\bottomrule
\end{longtable}

\begin{conceptbox}
These four overlap with the HEC 2023 elective \emph{Knowledge Representation and
Reasoning} (3 (3-0), prerequisite \emph{Artificial Intelligence}), and that is
intended rather than accidental. This course teaches each of them \emph{once, at
first pass}, because reasoning is not teachable without a representation to reason
over. The elective goes considerably further --- description logic, reasoning over
description logic, ontology languages, logical agents, commonsense reasoning and
explainable AI --- and none of that appears here. Appendix~F sets the line out in
full.
\end{conceptbox}

\section*{Three departures from the outline's grouping}

\begin{conceptbox}
The content of both outlines is taught in full. The \emph{grouping} departs in
three places, each deliberately.

\smallskip
\textbf{1. Search is taught before logic.} Both outlines list reasoning and
knowledge representation first. Here Part~II comes first, so that Chapters~7, 8
and~9 can reuse its machinery and so that students write working code in Week~3.

\smallskip
\textbf{2. Planning sits with search, not with knowledge.} STRIPS progression
search \emph{is} the frontier of Chapters~4 and~5 applied to a new state space.
Teaching it in Part~III keeps means--ends analysis next to the General Problem
Solver that uses it.

\smallskip
\textbf{3. Learning gets one signpost chapter, not a unit.} HEC 2017 names three
learning paradigms and HEC 2023 names artificial neural networks. All four clauses
are covered, in named sections of Chapter~17, and then the course stops: the
perceptron's exclusive-or limit is the last thing taught, and \emph{Machine
Learning} (Semester~6) takes over. Teaching more would duplicate that course and
leave less room for the symbolic material no other course covers.
\end{conceptbox}

\newpage

%=====================================================================
\chapter*{The Course at a Glance}
\addcontentsline{toc}{chapter}{The Course at a Glance}
\markboth{The Course at a Glance}{}

\dsfigh{%
\begin{tikzpicture}[
  part/.style={rounded corners=3pt, fill=#1, text=white, font=\bfseries\footnotesize,
               minimum width=2.5cm, text width=2.2cm, minimum height=0.95cm,
               align=center},
  ch/.style={rounded corners=2pt, draw=#1, line width=0.7pt, fill=#1!7, text=#1!60!black,
             font=\scriptsize, minimum width=2.5cm, minimum height=0.52cm, align=center},
  arr/.style={-{Stealth[length=2mm]}, primary!45, line width=0.8pt}
]
\node[part=primary]   (p1) at (0,0)    {Part I\\Foundations};
\node[part=secondary] (p2) at (3.0,0)  {Part II\\Search};
\node[part=greenhl]   (p3) at (6.0,0)  {Part III\\Constraints, Games, Plans};
\node[part=purplehl]  (p4) at (9.0,0)  {Part IV\\Knowledge and Logic};
\node[part=tealhl]    (p5) at (12.0,0) {Part V\\Uncertainty};
\node[part=goldhl]    (p6) at (15.0,0) {Part VI\\Learning and Practice};

\foreach \a/\b in {p1/p2,p2/p3,p3/p4,p4/p5,p5/p6}{\draw[arr] (\a) -- (\b);}

\foreach \i/\t in {1/1. What AI Is, 2/2. Agents and PEAS, 3/3. Representations}{
  \pgfmathsetmacro{\yy}{-0.90-(\i-1)*0.62}
  \node[ch=primary] at (0,\yy) {\t};}

\foreach \i/\t in {4/{}, 1/4. Uninformed, 2/5. Informed and A*, 3/6. Local Search}{
  \pgfmathsetmacro{\yy}{-0.90-(\i-1)*0.62}
  \node[ch=secondary] at (3.0,\yy) {\t};}

\foreach \i/\t in {1/7. Constraints, 2/8. Games, 3/9. Planning}{
  \pgfmathsetmacro{\yy}{-0.90-(\i-1)*0.62}
  \node[ch=greenhl] at (6.0,\yy) {\t};}

\foreach \i/\t in {1/10. Representation, 2/11. Propositional, 3/12. First-Order,
                   4/13. Expert Systems}{
  \pgfmathsetmacro{\yy}{-0.90-(\i-1)*0.62}
  \node[ch=purplehl] at (9.0,\yy) {\t};}

\foreach \i/\t in {1/14. Uncertainty, 2/15. Bayes Nets, 3/16. Fuzzy Logic}{
  \pgfmathsetmacro{\yy}{-0.90-(\i-1)*0.62}
  \node[ch=tealhl] at (12.0,\yy) {\t};}

\foreach \i/\t in {2/18. NLP, 3/19. Classic Systems, 4/20. Trends,
                   5/21. Responsible AI, 6/22. Project}{
  \pgfmathsetmacro{\yy}{-0.90-(\i-1)*0.62}
  \node[ch=goldhl] at (15.0,\yy) {\t};}
% the signpost chapter, marked
\node[rounded corners=2pt, draw=accent, line width=1.4pt, fill=accent!12,
      text=accent!60!black, font=\scriptsize\bfseries, minimum width=2.5cm,
      minimum height=0.52cm] at (15.0,-0.90) {17. Learning};

\node[rounded corners=3pt, fill=primary!7, draw=primary!30,
      minimum width=17.3cm, minimum height=0.66cm,
      font=\scriptsize\bfseries, text=primary] at (7.5,-4.55)
      {\faIcon{tasks}~~One running example throughout: IT project management ---
       task dependencies, sprints, tickets, releases, staffing and risk};
\end{tikzpicture}}%
{The six parts and the chapters in each. Chapter 17 (outlined) is the signpost to
\emph{Machine Learning}: the three learning paradigms and the perceptron, and no
further.}{coursemap}

\vspace{4pt}
{\small\itshape\color{gray!70!black}
This course is the published prerequisite for \emph{Machine Learning},
\emph{Programming for Artificial Intelligence}, \emph{Knowledge Representation and
Reasoning} and \emph{Introduction to Data Science}. Appendix~F states what it hands
to each of them. A one-page map of the whole
\emph{Artificial Intelligence $\rightarrow$ Machine Learning $\rightarrow$ Data
Science} progression is distributed alongside this handout.}

\newpage

%=====================================================================
\chapter*{The Running Example}
\addcontentsline{toc}{chapter}{The Running Example}
\markboth{The Running Example}{}

Every example in this book comes from IT project management. There are no other
application domains and no lens panels. The point is not that project management
is the most interesting use of AI; it is that a reader who has to learn a new
domain in every chapter is spending attention on the domain instead of on the
method. One domain, met repeatedly, becomes invisible --- which is what you want.

\rowcolors{2}{white}{lightbg}
\begin{longtable}{@{}L{2.1cm}L{5.3cm}L{7.8cm}@{}}
\toprule
\rowcolor{primary!15}
\textbf{Chapter} & \textbf{The recurring example} & \textbf{What it is used to
teach} \\
\midrule
\endhead
1 & Five tools in a project office & Which software is an AI system, and which is
merely useful \\
2 & Overnight ticket triage & PEAS, environment properties, and why an agent needs
internal state \\
3 & The nine-task release graph & States as values; the frontier; the
\texttt{Problem} interface \\
4 & Ordering the release tasks & Uninformed search, and the explored set \\
5 & Manual against automated testing & Heuristics, $A^{*}$, and why an investment
is invisible to greedy search \\
6 & Staffing four workstreams & Local maxima, ridges, annealing and evolution \\
7 & Sprint scheduling & Variables, domains, constraints, and propagation \\
8 & Negotiating a release slot & Minimax, alpha--beta, and move ordering \\
9 & Building a release plan & STRIPS operators and means--ends analysis \\
10 & The delivery ontology & Semantic networks, inheritance, frames and
defaults \\
11 & The release gate & Propositional inference three ways \\
12 & Who blocks whom & Quantifiers, unification and resolution \\
13 & The escalation adviser & Forward and backward chaining, and explanation \\
14 & Ship now or delay a week & Expected utility and the value of information \\
15 & Diagnosing a defect's cause & Bayesian network inference, four ways \\
16 & Rating project risk & Membership functions and Mamdani inference \\
17 & Predicting a late task & The perceptron, and the limit where this course
stops \\
18 & Ticket subject lines & Grammars, chart parsing and pattern substitution \\
19 & Four classic systems & Read through the machinery of Chapters 9, 12, 13
and 18 \\
20 & A tool-using delivery agent & What current systems add, and what they do not
change \\
21 & Auditing the escalation rules & Bias, explanation as evidence, and
computational limits \\
22 & Your own project & The whole course, end to end \\
\bottomrule
\end{longtable}

\newpage

%=====================================================================
\chapter*{Notation}
\addcontentsline{toc}{chapter}{Notation}
\markboth{Notation}{}

\rowcolors{2}{white}{lightbg}
\begin{longtable}{@{}L{2.8cm}L{12.4cm}@{}}
\toprule
\rowcolor{primary!15}
\textbf{Symbol} & \textbf{Meaning} \\
\midrule
\endhead
$s$, $s'$ & A state, and a successor of it \\
$b$, $d$, $m$ & Branching factor, depth of the shallowest solution, maximum depth
of the space \\
$g(n)$ & Cost of the path from the start to node $n$ \\
$h(n)$ & Heuristic estimate of the cost from $n$ to a goal \\
$h^{*}(n)$ & The true cheapest cost from $n$ to a goal \\
$f(n)$ & $g(n) + h(n)$, the estimate used by $A^{*}$ \\
$C^{*}$ & The cost of an optimal solution \\
$X_{i}$, $D_{i}$ & A constraint variable and its domain \\
$\alpha$, $\beta$ & The bounds carried by alpha--beta pruning \\
$\neg$, $\wedge$, $\vee$, $\Rightarrow$, $\Leftrightarrow$ & Not, and, or,
implies, if and only if \\
$\models$ & Entails: every model of the left is a model of the right \\
$\vdash$ & Derives: the right is provable from the left by the inference
procedure in use \\
$\forall$, $\exists$ & For all, there exists \\
$\theta$ & A substitution, as produced by unification \\
$P(a)$, $P(a \mid b)$ & Probability of $a$; probability of $a$ given $b$ \\
$\mathrm{EU}(a)$ & Expected utility of action $a$ \\
$\mu_{A}(x)$ & Degree of membership of $x$ in fuzzy set $A$ \\
$w_{i}$, $\eta$ & A perceptron weight and the learning rate \\
\bottomrule
\end{longtable}

\begin{alertbox}[One notational warning]
$\models$ and $\vdash$ look similar and mean different things.
$\mathit{KB} \models \alpha$ says $\alpha$ is \emph{true} in every world where the
knowledge base is true --- a fact about meaning. $\mathit{KB} \vdash \alpha$ says
$\alpha$ can be \emph{derived} by a particular procedure --- a fact about that
procedure. Chapter~11 shows why an inference procedure is worth having exactly
when the two coincide, and what it means when they do not.
\end{alertbox}

\newpage

%=====================================================================
\chapter*{Prerequisite Self-Check}
\addcontentsline{toc}{chapter}{Prerequisite Self-Check}
\markboth{Prerequisite Self-Check}{}

Ten questions. If you can answer eight, you are ready for Chapter~1. If you cannot,
the chapter named beside each one is where to look first --- and every one of them
is in Appendix~A.

\begin{enumerate}[leftmargin=2.4em, itemsep=4pt, topsep=4pt]
  \item What is the difference between a Python \texttt{list} and a
        \texttt{tuple}, and why can only one of them be a dictionary key?
        \hfill{\scriptsize\color{neutral}(\chref{python})}
  \item What does \texttt{\{1,2\} <= \{1,2,3\}} evaluate to, and what does the
        operator mean for sets? \hfill{\scriptsize\color{neutral}(\chref{python})}
  \item Write a class with one method that subclasses override. What is this
        pattern for? \hfill{\scriptsize\color{neutral}(Object Oriented
        Programming)}
  \item What does the \texttt{yield} keyword do, and how does a generator differ
        from a list? \hfill{\scriptsize\color{neutral}(\chref{python})}
  \item What is the time complexity of removing the first element of a Python
        list, and what should you use instead? \hfill{\scriptsize\color{neutral}(\chref{python})}
  \item Given a dictionary mapping each task to the set of tasks it depends on,
        write one expression giving the tasks with no unmet dependency.
        \hfill{\scriptsize\color{neutral}(\chref{python})}
  \item What is recursion, and what must every recursive function have?
        \hfill{\scriptsize\color{neutral}(Programming Fundamentals)}
  \item If a fair coin is tossed twice, what is the probability of at least one
        head? \hfill{\scriptsize\color{neutral}(\chref{uncertainty})}
  \item What does it mean for two events to be independent, and write the
        consequence as an equation. \hfill{\scriptsize\color{neutral}(\chref{uncertainty})}
  \item A truth table for two variables has how many rows? For $n$ variables?
        \hfill{\scriptsize\color{neutral}(\chref{propositional})}
\end{enumerate}

\begin{conceptbox}
Notice what is \emph{not} on that list: calculus, linear algebra, statistics and
any experience of machine learning. This course needs none of them. It needs you to
be fluent with sets, dictionaries, recursion and classes, and to be willing to work
an example by hand before trusting the code.
\end{conceptbox}

\newpage
\tableofcontents
\newpage
\listoffigures
\newpage

\pagenumbering{arabic}
